\section{Simultaneous removal of the lifted quadratic families}
\label{sec:simultaneous-quadratic-normalization}

The individual lifts can be normalized simultaneously, because their
arithmetic restrictions remain independent after the preceding
normalizations. We first prove the precise character comparison that
makes this assertion possible.

\begin{lemma}\label{lem:arithmetic-normalization-annihilator}
Let \(G=NH\) be a profinite group, where \(N\) is closed and normal
and \(H\) is closed. Put \(K=N\cap H\). Let \(R\subseteq N\)
be closed and normal in \(G\), with \(G=RH\), and put
\(J=R\cap H\). For any finite abelian group \(C\) with trivial
action, restriction induces an isomorphism
\[
 \ker\bigl(\operatorname{Hom}_{\rm cts}(N,C)^G
                 \to\operatorname{Hom}_{\rm cts}(R,C)\bigr)
 \xrightarrow{\ \sim\ }
 \ker\bigl(\operatorname{Hom}_{\rm cts}(K,C)^H
                 \to\operatorname{Hom}_{\rm cts}(J,C)\bigr).
\]
\end{lemma}
\begin{proof}
The factorization \(G=RH\) gives \(N=RK\): if \(n=rh\in N\)
and \(r\in R\subseteq N\), then \(h\in N\cap H=K\).
Consequently inclusion induces a topological isomorphism
\[
                         K/J\xrightarrow{\ \sim\ }N/R.
\]
It is continuous and bijective between compact Hausdorff groups.
An \(H\)-invariant character of \(K\) killing \(J\) therefore
defines a continuous character of \(N/R\). Conjugation by \(R\)
is trivial on this quotient, and \(G=RH\); thus \(H\)-invariance
is precisely the required \(G\)-invariance. Conversely every
\(G\)-invariant character of \(N\) killing \(R\) restricts to
one of the stated characters of \(K\). The quotient isomorphism
makes these constructions inverse.
\end{proof}

Return to \(G=\GT\), \(H=G_{\Q}\) and the preceding notation.
Let
\[
 \mathcal I=\{(\ell,q):\ell<q\text{ are odd primes}\},\qquad
 \mathcal P=\prod_{(\ell,q)\in\mathcal I}\F_2,
\]
and use the ordered lift \(\xi_{\ell,q}\) constructed from the
cyclic degree-\(\ell\) cover with coefficients modulo \(q\).
Write \(\mathfrak t_2\) for the entire coefficient-two cubic CM
homomorphism, with target
\[
 \mathcal U=\{z\in1+4\Z_2[\omega]:z\bar z=1\},
 \qquad \omega^2+\omega+1=0.
\]
Put \(\delta=1+2\omega\), and write
\[
 \tau=\frac{\log\mathfrak t_2}{4\delta}\pmod2,
 \qquad
 \Sigma_0=\F_2\eta_3\oplus
                 \bigoplus_{p\text{ odd}}\F_2\psi_p.
\]
Thus \(\tau\in X_2\), and its arithmetic restriction is
\(\kappa_{2,3}+\psi_3\), by
Proposition~\ref{prop:fr-cubic-two-Y2-class}.
Define the closed normal subgroups
\[
 \begin{split}
 R^*&=R_D\cap\ker\mathfrak t_2
                   \cap\bigcap_{(\ell,q)\in\mathcal I}\ker\xi_{\ell,q},\\
 V^*&=V_D\cap R^*.
 \end{split}
\]

\begin{theorem}\label{thm:simultaneous-quadratic-normalization}
The joint homomorphism
\[
 (\mathfrak t_2,(\xi_{\ell,q})_{\mathcal I})
                      :R_D\longrightarrow\mathcal U\times\mathcal P
\]
is surjective already on \(R_D\cap H\), and remains surjective
on \(V_D\cap H\). Consequently
\[
 R_D/R^*\simeq\mathcal U\times\mathcal P,
 \qquad V_D/V^*\simeq\mathcal U\times\mathcal P.
\]
Both quotient descriptions hold on the arithmetic intersections as
well. Moreover
\[
 G=R^*H,\qquad
 B\simeq R^*/\overline{[G,R^*](R^*\cap H)},
\]
and
\[
 B_{\max2}\simeq
 \left(V^*/\overline{[G,V^*](V^*\cap H)}\right)_{\max2}.
\]
Restriction from \(X_2\) to either \(R^*\) or \(V^*\) has kernel
\[
 \Sigma_0\oplus\F_2\tau\oplus
                 \bigoplus_{(\ell,q)\in\mathcal I}\F_2\xi_{\ell,q}.
\]
Arithmetic restriction is injective on this kernel. In particular
\(V^*\) is a proper subgroup of infinite index in \(V_D\),
already on the arithmetic intersection.
\end{theorem}
\begin{proof}
Put \(R'=R_D\cap\ker\mathfrak t_2\) and
\(V'=V_D\cap\ker\mathfrak t_2\). Theorem~\ref{thm:fr-cubic-two-exact-image}
and Corollary~\ref{cor:fr-cubic-two-joint-gate} give
\[
 G=R'H,\qquad R'/V'\text{ is pro-odd},\qquad
 \ker(X_2\to H^1(R',\F_2))=\Sigma_0\oplus\F_2\tau.
\]
The arithmetic restrictions of this last space are exactly the span
of \(\eta_3\), all \(\psi_p\), and \(\kappa_{2,3}\).
Theorem~\ref{thm:fr-Y2-complete} says that the characters
\(\kappa_{\ell,q}\), for \((\ell,q)\in\mathcal I\), are linearly
independent modulo that span. Applying
Lemma~\ref{lem:arithmetic-normalization-annihilator} with \(R=R'\)
shows that their restrictions to \(R'\cap H\) are linearly independent.
They are the arithmetic restrictions of the \(\xi_{\ell,q}\).

The image of the continuous map
\(R'\cap H\to\mathcal P\) is compact and hence closed. If it
were proper, some finite-coordinate quotient would also have proper
image: a point outside a closed subset has a basic open neighbourhood
disjoint from it. A proper vector subspace of that finite product
has a nonzero linear annihilator. This would give a nonzero finite
linear combination of the \(\kappa_{\ell,q}\) vanishing on
\(R'\cap H\), a contradiction. Therefore this map is onto.

The same independence holds on \(V'\cap H\). Indeed,
\[
 (R'\cap H)/(V'\cap H)\hookrightarrow R'/V'
\]
is pro-odd. A quadratic character vanishing on \(V'\cap H\)
would factor through this pro-odd quotient and so would vanish on
\(R'\cap H\). The preceding compact-image argument now gives
surjectivity \(V'\cap H\to\mathcal P\).

The coefficient-two theorem gives
\(\mathfrak t_2(R_D\cap H)=\mathfrak t_2(V_D\cap H)=\mathcal U\).
Given \((u,e)\in\mathcal U\times\mathcal P\), first choose an
arithmetic element with coefficient-two value \(u\). Multiplication
by an element of \(R'\cap H\), or of \(V'\cap H\), prescribes
its whole \(\mathcal P\)-coordinate without changing \(u\).
This proves both joint surjectivity assertions, including their full
group versions. Their kernels give the asserted quotient isomorphisms.

For \(r\in R_D\), choose \(h\in R_D\cap H\) with the same joint
image. Then \(rh^{-1}\in R^*\), so
\(R_D=R^*(R_D\cap H)\) and \(G=R^*H\).
The invariant-character extension argument for a closed normal
factor multiplying with \(H\) now gives the stated full relative
abelian quotient. More explicitly, for every finite abelian group
\(C\), a \(G\)-invariant homomorphism \(R^*\to C\) vanishing
on \(R^*\cap H\) extends uniquely to \(G\) by the rule
\(rh\mapsto\chi(r)\); compactness gives continuity. Taking all
finite targets identifies the profinite relative abelian quotients.

The natural injection
\(R^*/V^*\hookrightarrow R_D/V_D\) has pro-odd image.
Lemma~\ref{lem:relative-proodd-restriction}, applied with the proved
factorization \(G=R^*H\), identifies the maximal pro-two relative
quotient on \(V^*\).

Finally, a character in \(X_2\) vanishing on \(R^*\) restricts
to a continuous quadratic character of
\(R_D/R^*\simeq\mathcal U\times\mathcal P\).
Since \(\mathcal U\simeq\Z_2\), its only nonzero quadratic
character is \(\tau\). A continuous character of \(\mathcal P\)
depends on finitely many coordinates. Subtracting the corresponding
finite sum of \(\tau\) and the \(\xi_{\ell,q}\) leaves a character
vanishing on \(R_D\), hence one in \(\Sigma_0\).
Conversely all these characters vanish on \(R^*\).
The sum is direct, and has injective arithmetic restriction, by the
explicit basis of \(Y_2\). A character vanishing on \(V^*\)
vanishes on \(R^*\), since \(R^*/V^*\) is pro-odd. This proves
the identical kernel on \(V^*\) and completes the proof.
\end{proof}

This removes the cubic coefficient-two image and all odd-pair
characters simultaneously. The next construction lifts the remaining
dyadic-pair arithmetic characters; its application gives a further
normalization of the relative kernel.
